\documentclass{article}


% 中文支持
\usepackage[UTF8,heading = true]{ctex}


% 页面布局
\usepackage[a4paper]{geometry}
\usepackage{fancyhdr}
\pagestyle{fancy}
\newcommand{\Lhead}[1]{\fancyhead[L]{#1}}
\newcommand{\Chead}[1]{\fancyhead[C]{#1}}
\newcommand{\Rhead}[1]{\fancyhead[R]{#1}}
\newcommand{\Lfoot}[1]{\fancyfoot[L]{#1}}
\newcommand{\Cfoot}[1]{\fancyfoot[C]{#1}}
\newcommand{\Rfoot}[1]{\fancyfoot[R]{#1}}
\renewcommand{\headrulewidth}{0.4pt}
\renewcommand{\footrulewidth}{0pt}
\usepackage{lastpage}
%待填信息
% -----------------------------------------------------
\newcommand{\Course}{深度学习平台与应用}             		%      |
\newcommand{\Assignment}{深度学习作业 3}		  %      | 
\newcommand{\Name}{田永铭}				    	 %      |
\newcommand{\ID}{221900180}				  	   %       |
\newcommand{\Email}{tianym2022@smail.nju.edu.cn}  %    |
% ------------------------------------------------------
\fancyhead[L]{\Course}
\fancyhead[C]{\Assignment}
\fancyhead[R]{\ID \Name}
\fancyfoot[C]{\thepage\ / \pageref*{LastPage}}	
\setlength\headheight{36pt}
% 自定义页面样式第一页只有页脚没有页眉
\makeatletter
\def\ps@onlyfoot{
	\let\@oddhead\@empty
	\let\@evenhead\@empty
	\def\@oddfoot{\hfil\thepage\ / \pageref*{LastPage}\hfil}
	\def\@evenfoot{\hfil\thepage\ / \pageref*{LastPage}\hfil}
}
\makeatother
% maketitle
\title{\Huge\Assignment}
\author{\ID \Name \quad\href{mailto:tianym2022@smail.nju.edu.cn}{\color{magenta}\Email}}
\date{\today}


% 字体和颜色
\usepackage{fontspec} 
\fontsize{12pt}{\baselineskip}\selectfont
\usepackage{xcolor}


% 数学支持
\usepackage{amssymb}
\usepackage{amsfonts}
\usepackage{amsmath}
\usepackage{esint}
\usepackage{gensymb}
\usepackage{mathtools}
\usepackage{amsthm}

% 左边大括号
%\[
%\left\{
%\begin{aligned}
%	& \text{approximate  covariance matrix of GP} \\
%	& \text{only use subset}\\
%	& \text{conditional independence assumption}
%\end{aligned}
%\right.
%\]

% 矩阵
%\[
%\begin{matrix}
%	a & b \\
%	c & d
%\end{matrix}
%\]

% 行列式
%\[
%\begin{vmatrix}
%	a & b \\
%	c & d
%\end{vmatrix}
%\]


% 标题格式
\usepackage{titlesec}
% 自定义 section 标题格式
\renewcommand{\thesection}{\chinese{section}、}
\titleformat{\section}
{\normalfont\Large\bfseries\centering\color[rgb]{0,0,0.8}} % 设定字体和大小
{\thesection}{0em} % 设定编号和标题之间的距离
{}
% 自定义 subsection 标题格式
\renewcommand{\thesubsection}{\arabic{section}.\arabic{subsection} }
\titleformat{\subsection}
{\normalfont\large\bfseries\color[rgb]{0.23,0.65,0.99}}
{\thesubsection}{0em}
{}
% 自定义 subsubsection 标题格式
\renewcommand{\thesubsubsection}{\arabic{section}.\arabic{subsection}.\arabic{subsubsection} }
\titleformat{\subsubsection}
{\normalfont\normalsize\bfseries}
{\thesubsubsection}{0em}
{}
% 设置 section 标题的间距
\titlespacing*{\section}
{0pt}{3.5ex plus 1ex minus .2ex}{2.3ex plus .2ex}
% 设置 subsection 标题的间距
\titlespacing*{\subsection}
{0pt}{3.25ex plus 1ex minus .2ex}{1.5ex plus .2ex}
% 设置 subsubsection 标题的间距
\titlespacing*{\subsubsection}
{0pt}{3.25ex plus 1ex minus .2ex}{1.5ex plus .2ex}


% 代码 
% 如非默认，需使用[language=xxx,style=mystyle]
\usepackage{listings}
% language:  {[x86masm]Assembler}(汇编); python; C; C++; Java; SQL; Matlab; R; bash; make(makefile)
\lstdefinestyle{mystyle}{
	basicstyle=\ttfamily\small,
	numbers=left,
	numberstyle=\tiny\color{magenta},
	stepnumber=1,
	numbersep=10pt,
	backgroundcolor=\color[rgb]{0.95,0.95,0.95},
	showspaces=false,
	showstringspaces=false,
	showtabs=false,
	frame=single,
	rulecolor=\color{black},
	tabsize=2,
	captionpos=b,
	breaklines=true,
	breakatwhitespace=false,
	title=\lstname,
	keywordstyle=\color{blue},
	commentstyle=\color[rgb]{0.0,0.5,0.0},
	stringstyle=\color[rgb]{0.23,0.65,0.99}\itshape,
	escapeinside={\%*}{*)},
	morekeywords={}, % 添加更多关键字
	literate=
	{[}{{{\color{red}[}}}1
	{]}{{{\color{red}]}}}1
	{\{}{{{\color[rgb]{0.9,0,0.9}\{}}}1
	{\}}{{{\color[rgb]{0.9,0,0.9}\}}}}1
%	{0}{{{\color[rgb]{0.2,0.9,0.4}0}}}1
%	{1}{{{\color[rgb]{0.2,0.9,0.4}1}}}1
%	{2}{{{\color[rgb]{0.2,0.9,0.4}2}}}1
%	{3}{{{\color[rgb]{0.2,0.9,0.4}3}}}1
%	{4}{{{\color[rgb]{0.2,0.9,0.4}4}}}1
%	{5}{{{\color[rgb]{0.2,0.9,0.4}5}}}1
%	{6}{{{\color[rgb]{0.2,0.9,0.4}6}}}1
%	{7}{{{\color[rgb]{0.2,0.9,0.4}7}}}1
%	{8}{{{\color[rgb]{0.2,0.9,0.4}8}}}1
%	{9}{{{\color[rgb]{0.2,0.9,0.4}9}}}1
}
\lstset{style=mystyle}


% 图表
\usepackage{graphicx}
\usepackage{booktabs}
\usepackage{caption}
\captionsetup{font={small}}
\usepackage{diagbox}
\usepackage{array}
\usepackage{pgfplots}
\pgfplotsset{compat=1.17}

% 表1
%\begin{table}[h!]
%	\centering
%	\caption{Performance comparison on mini-ImageNet and tiered-ImageNet datasets.}
%	\begin{tabular}{@{}lcccccccc@{}}
%		\toprule
%		Dataset & Method & Backbone & 5-way 1-shot \\
%		\midrule
%		\multirow{9}{*}{\rotatebox[origin=c]{90}{mini-ImageNet}} 
%		& ProtoNet & Conv-4 & 46.291 \\
%		& MAML & Conv-4 & 49.704 \\
%		& BOIL & Conv-4 & 48.456 \\
%		& CAN & ResNet-12 & 59.493\\
%		& LEO & ResNet-12 & 55.056 \\
%		& SKD & ResNet-12 & 66.700  \\
%		& R2D2 & ResNet-12 & 59.802 \\
%		& DiffKendall (Ours) & ResNet-12 & \textbf{68.727} \\
%		\midrule
%		\multirow{6}{*}{\rotatebox[origin=c]{90}{tiered-ImageNet}} 
%		& ProtoNet & Conv-4 & 46.653 \\
%		& Relation Networks & Conv-4 & 55.107  \\
%		& MAML & Conv-4 & 51.038 \\
%		& R2D2 & ResNet-12 & 64.916  \\
%		& DiffKendall (Ours) & ResNet-12 & \textbf{66.598} \\
%		\bottomrule
%	\end{tabular}
%\end{table}

% 表2
%\begin{table}[t]
%	\caption{Classification accuracies for naive Bayes and flexible
%		Bayes on various data sets.}
%	\label{sample-table}
%	\vskip 0.15in
%	\begin{center}
%		\begin{small}
%			\begin{sc}
%				\begin{tabular}{lcccr}
%					\toprule
%					Data set & Naive & Flexible & Better? \\
%					\midrule
%					Breast    & 95.9$\pm$ 0.2& 96.7$\pm$ 0.2& $\surd$ \\
%					Cleveland & 83.3$\pm$ 0.6& 80.0$\pm$ 0.6& $\times$\\
%					Glass2    & 61.9$\pm$ 1.4& 83.8$\pm$ 0.7& $\surd$ \\
%					Credit    & 74.8$\pm$ 0.5& 78.3$\pm$ 0.6&         \\
%					Horse     & 73.3$\pm$ 0.9& 69.7$\pm$ 1.0& $\times$\\
%					Meta      & 67.1$\pm$ 0.6& 76.5$\pm$ 0.5& $\surd$ \\
%					Pima      & 75.1$\pm$ 0.6& 73.9$\pm$ 0.5&         \\
%					Vehicle   & 44.9$\pm$ 0.6& 61.5$\pm$ 0.4& $\surd$ \\
%					\bottomrule
%				\end{tabular}
%			\end{sc}
%		\end{small}
%	\end{center}
%	\vskip -0.1in
%\end{table}


% 算法
%\usepackage{algorithm}
%\usepackage{algorithmic}
%\begin{algorithm}[tb]
%	\caption{Bubble Sort}
%	\label{alg:example}
%	\begin{algorithmic}
	%		\STATE {\bfseries Input:} data $x_i$, size $m$
	%		\REPEAT
	%		\STATE Initialize $noChange = true$.
	%		\FOR{$i=1$ {\bfseries to} $m-1$}
	%		\IF{$x_i > x_{i+1}$}
	%		\STATE Swap $x_i$ and $x_{i+1}$
	%		\STATE $noChange = false$
	%		\ENDIF
	%		\ENDFOR
	%		\UNTIL{$noChange$ is $true$}
	%	\end{algorithmic}
%\end{algorithm}

% 超链接
\usepackage{hyperref}


% 论文参考文献
% 需在正文结尾加入 \bibliography{references}，在同目录下创建references.bib
\bibliographystyle{plain}


% 表情
% https://ctan.math.washington.edu/tex-archive/graphics/pgf/contrib/tikzsymbols/tikzsymbols.pdf
\usepackage{tikzsymbols}
\usepackage{tikz}
% 背景颜色：\colorbox{yellow}{...}
% 3D版本： 在前面加上d,比如：dSmiley

% 常用
% 笑脸：\Smiley[⟨scale⟩][⟨color⟩]
% 哭脸：\Sadey[⟨scale⟩][⟨color⟩]
% 生气脸： \Annoey[⟨scale⟩][⟨color⟩]
% 大笑： \Laughey[⟨scale⟩][⟨color⟩][⟨mouth color⟩]
% 挤眉弄眼： \Winkey[⟨scale⟩][⟨color⟩]
% 去世： \Xey[⟨scale⟩][⟨color⟩]
% 升天： \wInnocey[⟨scale⟩]
% 口吐芬芳： \Vomey[⟨scale⟩][⟨color⟩][⟨vomit color⟩]
% 猫： \Cat[⟨scale⟩]
% 炸弹： \Ninja[⟨scale⟩][⟨color⟩][⟨headband color⟩][⟨eye color⟩]
% 困： \Sleepey[⟨scale⟩][⟨color⟩][⟨cap color⟩][⟨star color⟩]
% 口罩：\Maskey[⟨scale⟩][⟨color⟩][⟨mask color⟩]
% 渐变弧度脸(-2到2)： \Changey[⟨scale⟩][⟨color⟩]{⟨mood⟩}
% 树（春夏秋冬）： \Springtree[⟨scale⟩]
% 咖啡： \Coffeecup[⟨scale⟩]



% 简化命令
\newcommand{\hs}[1]{\hspace{{#1}em}} %空格
\newcommand{\qd}{\quad} %空格
\newcommand{\qqd}{\qquad} %空格
% 链接
\newcommand{\myhref}[2]{\href{#1}{{\color[rgb]{0.9,0,0.9}{#2}}}}
\newcommand{\myred}[1]{{\color{red}{#1}}} %红色
% 图 args：figures下图片路径，大小，图片显示名
\newcommand{\myfig}[3]{
	\begin{figure}[h!]
		\centering
		\includegraphics[width=#2\linewidth]{figures/#1}
		\caption{#3}
		\label{fig:#1}
	\end{figure}
}
% 左右分栏插入图片
% args: 左图文件名，左图宽度，左图标题，右图文件名，右图宽度，右图标题
\newcommand{\mytwocolumnfig}[6]{
	\begin{figure}[h!]
		\centering
		\begin{minipage}[h!]{0.45\linewidth}
			\centering
			\includegraphics[width=#2\linewidth]{figures/#1}
			\caption{#3}
			\label{fig:#1}
		\end{minipage}
		\hfill
		\begin{minipage}[h!]{0.45\linewidth}
			\centering
			\includegraphics[width=#5\linewidth]{figures/#4}
			\caption{#6}
			\label{fig:#4}
		\end{minipage}
	\end{figure}
}


\newtheoremstyle{Problem}
{}
{}
{}
{}
{\bfseries}
{}
{ }
{Problem\thmnumber{ #2}\thmname{ #1}\thmnote{ (#3)} }
\theoremstyle{Problem}

\newtheorem{prob}{}
\newtheoremstyle{Solution}
{}
{}
{}
{}
{\bfseries}
{:}
{ }
{\thmname{#1}}
\makeatletter
\makeatother
\theoremstyle{Solution}
\newtheorem*{sol}{Solution}


\begin{document}
	\maketitle
	
	\section{视频理解}
	\begin{prob}
		\textbf{分别用简短的语言描述一下单帧 CNN 模型，Early Fusion，Late Fusion，3D CNN的做法以及区别。}
	\end{prob}
	\begin{sol}
		{\color{red} 做法：}
		
		\textbf{单帧 CNN 模型}：  
		训练普通的 2D CNN 独立对视频单一图象帧进行分类。
		
		\textbf{Early Fusion}： 第一个2D卷积处理所有时间信息，网络的其他部分是常规的2D卷积。
		
		\textbf{Late Fusion}： 将每帧的CNN特征并联（或者池化）并送入 MLP。		

		\textbf{3D CNN}：  扩展传统的2D卷积神经网络，加入时序维度T，使用三维卷积核来处理时序数据或视频帧，能够捕捉时空信息的变化。
		
		{\color{red} 区别：}  

			\textbf{单帧 CNN} 处理的是静态图像，是最一般，最难捕捉动态信息的方法。而\textbf{3D CNN}由于用了三维卷积，是最适合能够处理时序或视频数据，能捕捉动态信息的。
			
			\textbf{其它两种方法}介于前两者之间，都是用``融合"的思想，但是\textbf{Early Fusion} 是将不同模态的原始数据融合在一起输入网络，而 \textbf{Late Fusion} 则是在每个模态单独处理后再进行融合。\\
	\end{sol}
	
	
	\begin{prob}
		\textbf{如何使用RNN（如LSTM）建模视频中的长时间依赖？相比于3D CNN，这种方法有何特点？}
	\end{prob}
	
	\begin{sol}
	{\color{red} 使用 LSTM 进行视频长时间依赖建模:}
	  
	通常将视频切分成连续的帧序列，每一帧通过卷积神经网络（如 ResNet）提取特征，之后将这些特征作为 LSTM 网络的输入，由LSTM处理局部特征。基本流程见图~\ref{fig:1}。
	
	而具体用LSTM建模的时候，会考虑更多问题，比如有时为了节省显存，不训练 CNN预训练并将其用作特征提取器；又如使用多层 RNN来综合局部时间结构和全局时间结构等等。
	
\begin{figure}[h!]
	\centering
	\includegraphics[width=1\linewidth]{Figures/1}
	\caption{LSTM建模长时间序列}
	\label{fig:1}
\end{figure}

	\myred{相比3D CNN的特点:}
	
	\begin{itemize}
		\item \textbf{建模方式}：  
		\textbf{LSTM} 逐帧处理视频中的时序信息，通过内部的记忆单元来捕捉长时间的依赖关系，主要解决的是时序信息，空间信息还是由原本的CNN来完成。而 \textbf{3D CNN} 通过三维卷积直接对视频的时空数据一起处理，能够同时在空间和时间维度上进行特征提取。
		
		\item \textbf{捕捉时间依赖的能力}：  
		\textbf{LSTM} 由于由记忆单元等等，相比于\textbf{3D CNN}更擅长建模更长时间依赖，能够处理视频中跨多个时间步的长期信息（如情节变化、长时间动作）。
		
		\item \textbf{灵活度}：
		\textbf{LSTM} 可处理变长输入（任意帧数视频）,而\textbf{3D CNN}通常需要固定输入长度（帧数固定）。
		
		\item \textbf{使用场景}：
		\textbf{LSTM} 相比 \textbf{3D CNN}更适合处理长序列。
		
		\item \textbf{计算效率}：  
		\textbf{LSTM} 在处理长时间序列时计算成本较高，尤其是在长视频中，训练和推理可能变得非常耗时（要计算很多``门"）。而 \textbf{3D CNN} 可以通过多维卷积核并行处理时间帧，通常在处理大规模视频数据时具有更高的计算效率。但是它参数量通常比较大，需要更大的计算资源和内存。\\
	\end{itemize}
	
	\end{sol}
	
	\begin{prob}
		\textbf{结合课件中Spatio-Temporal Self-Attention (Nonlocal Block)部分的介绍，请补充下面的代码。}
		\begin{lstlisting}[language=python,style=mystyle]
import torch
import torch.nn as nn
class NonLocalBlock(nn.Module):
	def __init__(self, in_channels):
		super(NonLocalBlock, self).__init__()
		self.in_channels = in_channels
		self.inter_channels = in_channels // 2  # 缩减通道数，节约计算
		# 1x1x1 卷积层用于生成 Query, Key, Value
		self.query_conv = nn.Conv3d(in_channels, self.inter_channels, 
		kernel_size=1)
		self.key_conv = nn.Conv3d(in_channels, self.inter_channels, 
		kernel_size=1)
		self.value_conv = nn.Conv3d(in_channels, self.inter_channels, 
		kernel_size=1)
		# 最终的输出通道映射回原始的 in_channels
		self.out_conv = nn.Conv3d(self.inter_channels, in_channels, 
		kernel_size=1)
		# 用于归一化的 softmax
		self.softmax = nn.Softmax(dim=-1)
	def forward(self, x):
		# 输入维度: (N, C, T, H, W)
		batch_size, C, T, H, W = x.size()
		# 生成 Query, Key, Value
		query = self.query_conv(x).view(batch_size, self.inter_channels, -1)  # (N, C', T*H*W)
		key = self.key_conv(x).view(batch_size, self.inter_channels, -1)  #  (N, C', T*H*W)
		value = self.value_conv(x).view(batch_size, self.inter_channels, -1) # (N, C', T*H*W)
		# 请在下面补充完整 Spatio-Temporal Self-Attention 计算过程
		return out
# 测试模块
if __name__ == "__main__":
# 输入示例: Batch = 2, Channels = 64, Time = 8, Height = 32, Width = 32
	x = torch.rand(2, 64, 8, 32, 32)
	nonlocal_block = NonLocalBlock(in_channels=64)
	out = nonlocal_block(x)
	print("输入维度:", x.shape)
	print("输出维度:", out.shape)
		\end{lstlisting}
	\end{prob}
	
	\newpage
	
	\begin{sol}
		补充部分的代码如下，主要采用的操作是变换形状和矩阵相乘。
\begin{lstlisting}[language=python,style=mystyle]
# 计算 QK^T
attn = torch.bmm(query.permute(0, 2, 1), key)  # (N, T*H*W, T*H*W)
# 归一化
attn = self.softmax(attn)  # (N, T*H*W, T*H*W)
# 计算加权值
out = torch.bmm(attn, value.permute(0, 2, 1))  # (N, T*H*W, C')
# 恢复原始维度
out = out.view(batch_size, self.inter_channels, T, H, W)
# 映射回原始通道数
out = self.out_conv(out)  # (N, C, T, H, W)
out = out + x  # 残差连接 
return out
\end{lstlisting}

生成结果符合预期,输出如图~\ref{fig:2}，函数作用前后维度不变。
\begin{figure}[h!]
	\centering
	\includegraphics[width=1\linewidth]{Figures/2}
	\caption{输入输出维度}
	\label{fig:2}
\end{figure}

	\end{sol}
		
	\section{风格迁移与可视化}
		
	\begin{prob}
		\textbf{ 在风格迁移中，Gram矩阵的作用是什么？}
	\end{prob}
	
	\begin{sol}
	Gram矩阵是图像的 HW 个向量之间两两的内积所组成的矩阵,反映了特征图之间的全局统计相关性，这种相关性与图像的纹理和风格密切相关,所以用于捕捉图像的“风格”信息。具体做法如下：
	
	\textbf{1. 风格表示：}
	
	图像的风格主要由颜色、纹理、形状等元素的统计特征决定。Gram矩阵通过计算图像中各个特征图（feature maps）之间的相关性来表示这种风格特征。
	
	设图像经过卷积神经网络（CNN）处理后得到一系列特征图（如在某层的输出），Gram矩阵就是这些特征图的内积矩阵，用于捕捉特征图之间的空间相关性。
	
	\textbf{2. 计算方法：}
	
	假设某一层的特征图大小为 \(C \times H \times W\)，其中 \(C\) 是特征图的数量（通道数），\(H\) 和 \(W\) 是特征图的高度和宽度。
	
	对于每个特征图 \(f_i\)（在二维图像中，通常是特征图的矩阵），Gram矩阵 \(G\) 的计算公式为：
	
	\[
	G_{ij} = \sum_{h=1}^{H} \sum_{w=1}^{W} f_i(h, w) f_j(h, w)
	\]
	
	即，Gram矩阵的每一个元素表示特征图 \(i\) 和特征图 \(j\) 的相关性。
	
	\textbf{3. 风格迁移：}
	
	在风格迁移的过程中，我们希望保持目标图像的内容结构，同时使其风格接近风格图像。因此，风格损失（style loss）是通过计算生成图像和风格图像的Gram矩阵之间的差异来度量风格相似性的。
	
	损失函数的一部分通常是：
	
	\[
	\text{Style Loss} = \sum_{l} \lambda_l \left( \| G_l^{\text{target}} - G_l^{\text{style}} \|^2 \right)
	\]
	
	其中，\(G_l^{\text{target}}\) 和 \(G_l^{\text{style}}\) 分别是目标图像和风格图像在某层的Gram矩阵，\(l\) 表示不同层的贡献，\(\lambda_l\) 是该层的权重。
	
	通过这种方式，Gram矩阵使得风格迁移不仅关注图像的局部细节，还能够捕捉图像全局的风格特征。
	\end{sol}
	
	
	\begin{prob}
		\textbf{ 风格迁移和快速风格迁移的区别？}
	\end{prob}
	
	\begin{sol}
		思想上，从\textbf{风格迁移}到\textbf{快速风格迁移}，有点类似``不再处理单一像素，而是从特征出发"。具体来说，\textbf{风格迁移}需要通过 VGG 进行多次向前/向后传播，速度很慢。\textbf{快速风格迁移}使用前馈网络进行风格迁移，使用预训练的 CNN 计算与以前相同的损失,训练后使用单次前向传播对图像进行风格迁移，速度快。
		
		通过表\ref{table}可以直观看出区别：
		\begin{table}[h!]
			\centering
			\begin{tabular}{|c|c|c|}
				\hline
				\textbf{特性} & \textbf{风格迁移（Style Transfer）} & \textbf{快速风格迁移（Fast Style Transfer）} \\
				\hline
				\textbf{方法} & 基于优化过程，通过迭代更新生成图像 & 使用训练好的神经网络进行前馈计算 \\
				\hline
				\textbf{计算时间} & 计算时间较长，通常需要几百到几千次迭代 & 生成速度快，通常可以实时生成图像 \\
				\hline
				\textbf{生成速度} & 慢 & 快（接近实时） \\
				\hline
				\textbf{图像质量} & 较高，细节丰富，风格表现较好 & 可能有所牺牲，但仍能呈现较好的风格 \\
				\hline
				\textbf{实现方式} & 需要每次进行优化和计算损失 & 通过训练一个生成网络，一次训练后即可使用 \\
				\hline
				\textbf{计算开销} & 高，需进行多次前向传播和反向传播 & 低，一次前向传播即可生成图像 \\
				\hline
			\end{tabular}
			\caption{风格迁移与快速风格迁移的比较}
			\label{table}
		\end{table}
	\end{sol}
	

	\begin{prob}
		\textbf{显著性图是怎么计算的（Saliency Map），它的用途是什么？}
	\end{prob}
	
	\begin{sol}
	\myred{计算方法：}
	通过梯度反向传播找到重要的像素。首先进行前向传播，然后计算每个类预测值相对于图像像素的梯度，在 RGB 通道上取绝对值的最大值，具体计算方法如下：
	
	\begin{itemize}
		\item \textbf{Grad-CAM（梯度加权类激活映射）}: 通过计算类别分数对最后卷积层输出的梯度，生成显著性图：
		\[
		\text{Grad-CAM}(x) = \text{ReLU}\left( \sum_{k} \alpha_k A^k \right)
		\]
		其中，\( A^k \) 为卷积层特征图，\(\alpha_k\) 为加权系数。
		
		\item \textbf{CAM（类激活映射）}: 直接加权卷积层特征图的每个通道，生成显著性图，适用于全局平均池化后的网络结构。
		
		\item \textbf{基本梯度显著性图}: 通过类别分数对图像像素的梯度计算显著性图：
		\[
		S(x) = \left| \frac{\partial f(x)}{\partial x} \right|
		\]
		其中，\( f(x) \) 是网络的输出。
	\end{itemize}
	\myred{用途：}
	\begin{itemize}
		\item \textbf{可以用于无监督分割}：可以根据显著性划分区域。
		\item \textbf{可以用于发现数据偏差}:比如发现二哈被分类为狼的原因是因为数据图中有雪。
		\item \textbf{模型可解释性}: 帮助理解模型做出预测的依据。
		\item \textbf{特征可视化}: 可视化模型在不同层次学习的特征。
		\item \textbf{数据增强}: 通过关注显著性区域进行数据增强，提高模型性能。
		\item \textbf{图像编辑}: 用于图像修复、生成等任务，聚焦重要区域进行操作。\\
	\end{itemize}
	\end{sol}
	
	\begin{prob}
		\textbf{请补充完整下面显著性可视化的代码。}
		\begin{lstlisting}[language=python,style=mystyle]
import torch
import torchvision
import torchvision.transforms as T
import numpy as np
import matplotlib.pyplot as plt
import requests
from PIL import Image

def download(url,fname):
	response = requests.get(url)
	with open(fname,"wb") as f:
	f.write(response.content)
	
def preprocess(image, size=224):
	transform = T.Compose([
	T.Resize((size,size)),
	T.ToTensor(),
	T.Normalize(mean=[0.485, 0.456, 0.406], std=[0.229, 0.224, 0.225]),
	T.Lambda(lambda x: x[None]),
	])
	return transform(image)
	
def deprocess(image):
	transform = T.Compose([
	T.Lambda(lambda x: x[0]),
	T.Normalize(mean=[0, 0, 0], std=[4.3668, 4.4643, 4.4444]),
	T.Normalize(mean=[-0.485, -0.456, -0.406], std=[1, 1, 1]),
	T.ToPILImage(),
	])
	return transform(image)
	
def show_img(PIL_IMG):
	plt.imshow(np.asarray(PIL_IMG))
	
if __name__ =='__main__':
	model = torchvision.models.vgg19(pretrained=True)
	for param in model.parameters():
	param.requires_grad = False

	download("https://bkimg.cdn.bcebos.com/pic/3bf33a87e950352ac65cae81db13ecf2b21192131da3?x-bce-process=image/format,f_auto/quality,Q_70/resize,m_lfit,limit_1,w_536","input.jpg")
	img = Image.open('input.jpg') # 这里可以替换为自己的图片 
	X = preprocess(img) # X.shape: 1, 3, 224, 224
	model.eval()
	X.requires_grad_()
	output = model(X) # 1, 1000 
	print(output.shape)
	score_max_index = output.argmax()
	score_max = output[0, score_max_index]

#####
# 需要实现课件上显著性可视化部分，具体的步骤为：
# a) 通过反向传播获取梯度
# b) 在RGB通道上取梯度绝对值的最大值
# 请在下面补充代码：
#####

	plt.imshow(saliency[0], cmap=plt.cm.hot)
	plt.axis('off')
	plt.savefig("output.png")
		\end{lstlisting}
	\end{prob}
	
	\begin{sol}
	补充部分代码如下：
	\begin{lstlisting}[language=python,style=mystyle]
# a) 通过反向传播获取梯度
score_max.backward()  # 对最高得分类别反向传播
gradients = X.grad.data  # 获取输入的梯度

# b) 在 RGB 通道上取梯度绝对值的最大值
saliency, _ = torch.max(gradients.abs(), dim=1)  # 取每个像素点在 RGB 通道的最大绝对值
saliency = saliency.squeeze().cpu().numpy()  # 转为 numpy 格式，去掉 batch 维
	\end{lstlisting}
	生成结果如图~\ref{fig:3}。有点原图奶龙的亚子力。\Winkey[2][yellow]
	
\begin{figure}[h!]
	\centering
	\includegraphics[width=0.4\linewidth]{Figures/3}
	\caption{显著性图}
	\label{fig:3}
\end{figure}
	
	\end{sol}
	

\end{document}
